% !TeX root = ../教材题库_全部条目_转写工作稿.tex
% Source: XJ-B1-C05-HW-P202-Y0822
% 第5章 / 5.5 三角函数模型的简单应用 / 教材 p.202 / PDF 第209页
% OCR未校对：公式与文字均可能有误；未完成数学TeX转写；禁止直接发给学生
% 原题及相邻原解截图（核对用）：../题目原图/XJ-B1-C05-HW-P202-Y0822-01.webp
% 本文件只在转写完成后作为题干使用；OCR 条目尚不是可靠题干。
\begin{bankitem}{习题5.5 第4题}{第5章 / 5.5 三角函数模型的简单应用 / 教材 p.202 / PDF 第209页}{XJ-B1-C05-HW-P202-Y0822}{OCR未校对：公式与文字均可能有误；未完成数学TeX转写；禁止直接发给学生}
\begin{ocrdraft}
[教材p.202 / PDF第209页]
\par
4．如图，某地夏天8一14时的用电量变化曲线近似
\par
Y/（\ensuremath{\times}104kW\ensuremath{\cdot}h)
\par
满足函数y==Asin(oxr十@十乃（0<<丌）．
\par
50
\par
40“
\par
（1）求这一天的最大用电量及最小用电量；
\par
（2）写出这段曲线的函数解析式．
\par
30
\end{ocrdraft}
\end{bankitem}
